\documentclass[a4paper]{article}

\usepackage{graphicx}
\usepackage{amsmath,amssymb}
\usepackage{minted}
\usepackage{multirow}
\usepackage{float}
\usepackage{todonotes}
\usepackage{forest}
\usepackage{xstring}
\usepackage{xcolor}
\usepackage[most]{tcolorbox}
\usepackage{xparse}
\usepackage{natbib}

\usetikzlibrary{graphs,graphdrawing}
\usegdlibrary{layered}

% for external graphviz
\input{/home/donkarlo/Dropbox/repo/nd_ai_project/src/nd_ai/knoweldge_representation/language/formal/computer/programming/tex/latex/code_clips/external_graphviz.tex}

% for tags
\input{/home/donkarlo/Dropbox/repo/nd_ai_project/src/nd_ai/knoweldge_representation/language/formal/computer/programming/tex/latex/parts/control_sequence/kind/semtag/semtag.tex}

% for links
\input{/home/donkarlo/Dropbox/repo/nd_ai_project/src/nd_ai/knoweldge_representation/language/formal/computer/programming/tex/latex/code_clips/seehere.tex}

\title{Language}
\date{}

\begin{document}
    \maketitle
    All branches in wikipedia
    \url{https://en.wikipedia.org/wiki/Portal:Linguistics}
    
    
    \section{Definition from Saussure prespective}
        Language is a symbolic system that, in its most basic form, consists of sequences of replaceable components used to convey different meanings \cite{saussure-1916-course-in-general-linguistics}
        
        
    \section{Definition}
        Language is a system of conventional symbols (spoken, written, or signed) that enables members of a community to communicate ideas, experiences, and intentions.
        Encyclopedia Britannica
        
        \includegraphics[width=1\linewidth]{\figspath language_definition.png}
        
        \subsection{Phonetics}
        
        \subsection{Phonology}
        
        \subsection{Morphology}
        
        \subsection{Syntax}
        
        \subsection{Semantics}
        
        \subsection{Pragmatics}
    
    
    \section{Applied Linguistics}
    
    
    \section{Interdiciplinary}
        
        \subsection{Psycholinguistics}
        
        \subsection{Sociolinguistics}
        
        \subsection{Historical Linguistics}
        
        \subsection{Anthropological Linguistics}
        
        \subsection{Computational Linguistics}
        
        \subsection{Cognitive Linguistics}
        
        \subsection{Corpus Linguistics}
        
        \subsection{Biolinguistics}
            How does language ability roots in body and genetics
            Noam Chomsky \"Language Faculty\"
        
        \subsection{Neurolinguistics}
    
    
    \section{Cognitive Neuroscience of Language}
    
    
    \section{Language Acquisition}
        \url{https://en.wikipedia.org/wiki/Language_acquisition}
        \bibliographystyle{plainnat}
        \bibliography{/home/donkarlo/Dropbox/repo/nd_shared_research/bibliography/bibliography.bib}
\end{document}
